\section{Problem Formulation and Protocol Value}
\label{sec:formulation}
Consider a benchmark collected under one observation protocol and a future study
that may use another. The benchmark contains the measurements produced by the
realised protocol together with a trajectory-level target used for prediction.
The full latent trajectory need not be observed. An alternative protocol may
include measurements that were not collected in the benchmark. We first define
the target and protocol measurements in continuous time and then introduce the
discrete Gaussian model used in the theoretical analysis.

An \textit{observation protocol} is a finite collection of acquisition actions. An \textit{observation design} is the decision problem of choosing a protocol from a feasible family. The \textit{realised protocol} is the one that generated the benchmark data; an \textit{alternative protocol} is a feasible protocol whose predictive value is to be evaluated even though the benchmark does not contain the complete measurement vector that it would produce.

\subsection{Temporal Aggregate Targets and Observation Protocols}
\label{sec:temporal-protocol}

For unit $i$, let $Z_i=\{Z_i(t):t\in[0,T]\}$ be a latent scalar trajectory. The
supervised target is one number obtained from the whole trajectory. Specifically,
we first transform the state at each time by a function $g$ and then average the
transformed values over the observation horizon:
\begin{equation}
\Theta_{i,g}=\int_0^T\omega_T(t)\,g\{Z_i(t)\}\dd t,
\qquad \omega_T(t)\ge0,
\qquad \int_0^T\omega_T(t)\dd t=1 .
\label{eq:label-continuous}
\end{equation}
The weight function $\omega_T$ specifies which parts of the horizon contribute
to the target. Uniform weights, $\omega_T(t)=1/T$, give an ordinary time
average; non-uniform weights can emphasise particular periods.

Two choices of $g$ recur throughout the paper. If $g(z)=z$, then
$\Theta_{i,g}$ is the average level of the latent process, such as average
electricity demand over a day. If
$g_c(z)=\Ind\{z>c\}$, then $\Theta_{i,g_c}$ is the weighted fraction of time for
which the process exceeds $c$, and therefore represents exceedance-time burden.
Sleep-stage proportions, atrial-fibrillation burden and time spent in a
symptomatic state are examples of occupation-type targets.
The target is fixed when protocols are compared: changing the protocol changes
what is measured, not the quantity to be predicted.

One acquisition action $a$ returns a noisy linear measurement of the
trajectory,
$Y_{i,a}=L_aZ_i+\varepsilon_{i,a}$, where $L_a$ may represent a point
measurement or an average over a time window. We assume
$\varepsilon_{i,a}\sim\cN(0,r_a)$, independently across actions and units and
independently of $Z_i$. Each action has cost $c_a$. For a cost budget
$\mathcal B>0$, a protocol is a finite set $S=\{a_1,\ldots,a_D\}$ satisfying
$\sum_{a\in S}c_a\le\mathcal B$, and $\Pi_{\mathcal B}$ denotes the collection
of all such feasible protocols. Identification and estimation use this action
representation; \cref{sec:design} specifies how actions encode timing, support,
repetition, noise and cost.

\subsection{Discrete Gaussian Model}
\label{sec:discrete}

The theoretical analysis uses a grid $t_1<\cdots<t_p$ on $[0,T]$. We collect
the latent states on this grid in
$Z_i=(Z_i(t_1),\ldots,Z_i(t_p))^\top$ and assume
\begin{equation}
Z_i\sim\cN(0,K),\qquad \diag(K)=\Ind_p.
\label{eq:discrete-model}
\end{equation}
Thus $K$ records the dependence between times, while every marginal state has
mean zero and variance one. This standardisation fixes the scale on which $g$
is defined. Here $\diag(M)$ is the vector of diagonal entries of $M$, whereas
$\Diag(v)$ is the diagonal matrix with diagonal $v$.

Let $\omega=(\omega_1,\ldots,\omega_p)^\top$ be non-negative quadrature weights
with $\Ind_p^\top\omega=1$. The continuous target and the protocol measurements
then become
\begin{equation}
\Theta_{i,g}=\sum_{j=1}^p\omega_j\,g(Z_{ij}),
\qquad
Y_{i,S}=A_SZ_i+\varepsilon_{i,S},
\qquad
\varepsilon_{i,S}\sim\cN(0,R_S).
\label{eq:discrete-protocol}
\end{equation}
Each row $\ell_a^\top$ of $A_S\in\R^{D\times p}$ describes one action. A point
measurement at $t_j$ has $\ell_a=e_j$. A window measurement has non-negative
entries on the grid points covered by that window and is normalised so that
$\Ind_p^\top\ell_a=1$. With independent action noise,
$R_S=\Diag(r_a)_{a\in S}$. In short, $K$ describes the latent trajectory,
$\omega$ describes the target, and $(A_S,R_S)$ describes what protocol $S$
observes.

The empty protocol has no observations and explains no target variation. We
set
\begin{equation}
Q_\varnothing(K)=0_{p\times p},\qquad \cI_g(\varnothing;K)=0.
\label{eq:empty-protocol}
\end{equation}

\subsection{Protocol Value under the Gaussian Model}
\label{sec:protocol-values}

We value a protocol by the fraction of target variance explained by the best
population predictor based on its measurements.

\begin{definition}[Bayes value of an observation protocol]
\label{def:values}
Let $\Theta\in L^2$ be a scalar target with $\Var(\Theta)>0$. The Bayes protocol
value under squared loss is
\begin{equation}
\cI_{\mathrm{Bayes}}(S)=\frac{\Var\{\E(\Theta\mid Y_S)\}}{\Var(\Theta)}.
\label{eq:bayes-value}
\end{equation}
It is the population $R^2$ of the optimal predictor based on $Y_S$ and lies in
$[0,1]$.
\end{definition}

The best-linear analogue restricts the predictor to be affine. With
$\Sigma_S=\Var(Y_S)$ and $c_S=\Cov(Y_S,\Theta)$, it is
\begin{equation}
\cI_{\mathrm L}(S)=\frac{c_S^\top\Sigma_S^{+}c_S}{\Var(\Theta)},
\label{eq:linear-value}
\end{equation}
where $\Sigma_S^{+}$ is the Moore--Penrose inverse
\citep{penrose1955generalized}.
It depends only on second moments and satisfies
$\cI_{\mathrm L}(S)\le\cI_{\mathrm{Bayes}}(S)$, with equality when
$\E(\Theta\mid Y_S)$ is affine in $Y_S$.

Equations~\eqref{eq:bayes-value} and \eqref{eq:linear-value} define population
quantities. The empirical analyses evaluate fitted predictors by pooled
held-out cross-fitted $R^2$.

Under the discrete Gaussian model, the Bayes value can be calculated from $K$,
$g$ and the protocol matrices. The target transformation $g$ converts a latent
Gaussian correlation into a covariance as follows. For standard bivariate
normal variables
$(U,V_r)$ with correlation $r$, define
\begin{equation}
C_g(r)=\Cov\{g(U),g(V_r)\}
=\sum_{k\ge1}\frac{a_k^2}{k!}\,r^k,
\qquad g\in L^2(\phi),
\label{eq:Cg}
\end{equation}
where $\phi$ is the standard normal distribution and
$g-\E g=\sum_{k\ge1}(a_k/k!)H_k$ is the Hermite expansion of $g$ in
$L^2(\phi)$. In particular,
$\Cov\{g(Z_j),g(Z_k)\}=C_g(K_{jk})$. All dependence on the possibly nonlinear
target function enters the value calculation through $C_g$.

The protocol enters through the part of the latent trajectory recoverable from
its measurements. Let $\Sigma_S(K)=A_SKA_S^\top+R_S$. For a non-empty protocol
with nonsingular $\Sigma_S(K)$, Gaussian conditioning gives
\begin{equation}
Q_S(K)=KA_S^\top\Sigma_S(K)^{-1}A_SK
      =\Cov\{\E(Z\mid Y_S)\}.
\label{eq:QP}
\end{equation}
Thus $Q_S(K)$ is the covariance of the part of the trajectory recovered from
$Y_S$. The residual covariance used in the marginal-gain calculation is
introduced in \cref{sec:design}.

Combining this covariance transform with Gaussian conditioning yields the
protocol-value identity below \citep{zhang2026label}.

\begin{proposition}[Gaussian protocol-value identity; Zhang, 2026]
\label{prop:risk-identity}
Under \eqref{eq:discrete-model}--\eqref{eq:discrete-protocol}, let
$g\in L^2(\phi)$ and assume $V_g(K)>0$. Define
\begin{equation}
V_g(K)=\sum_{j,k=1}^p\omega_j\omega_k\,C_g(K_{jk}),
\qquad
F_g(S;K)=\sum_{j,k=1}^p\omega_j\omega_k\,C_g\{Q_S(K)_{jk}\}.
\label{eq:VF}
\end{equation}
Then $V_g(K)=\Var(\Theta_g)$ and
$F_g(S;K)=\Var\{\E(\Theta_g\mid Y_S)\}$, and
\begin{equation}
\cI_g(S;K)=\cI_{\mathrm{Bayes}}(S)=\frac{F_g(S;K)}{V_g(K)}.
\label{eq:ceiling}
\end{equation}
\end{proposition}

Thus $V_g(K)$ is total target variance, $F_g(S;K)$ is the part explained by
protocol $S$, and $\cI_g(S;K)$ is the corresponding Bayes-optimal population
$R^2$ based on $Y_S$. A proof is given in
Appendix~\ref{sec:app-identification}.

For two target functions, the covariance transform has a simple form. For the
mean target $g(z)=z$, $C_g(r)=r$, so
$V_g(K)=\omega^\top K\omega$ and
$F_g(S;K)=\omega^\top Q_S(K)\omega$; equivalently,
$\Theta=\omega^\top Z$. For the occupation target
$g_c(z)=\Ind\{z>c\}$, Plackett's identity \citep{plackett1954reduction} gives
\begin{equation}
C_{g_c}(r)=G_c(r)
=\int_0^r\frac{\exp\{-c^2/(1+s)\}}{2\pi\sqrt{1-s^2}}\dd s,
\label{eq:plackett}
\end{equation}
with $G_0(r)=\arcsin(r)/(2\pi)$.
